Els models de processament d'àudio basats en xarxes neuronals profundes (Deep Neural Networks, \acrshort{DNN}) sovint no generalitzen bé en condicions reals, on el rendiment es pot degradar substancialment. Aquesta tesi pretén reduir aquesta bretxa millorant la validesa ecològica de les dades (sovint sintètiques), en tres tasques de processament d'àudio diferents:

\begin{itemize}

\item \textbf{\textit{(i)}} Per a la millora de la parla (Speech Enhancement, \acrshort{SE}) binaural en audiòfons (Hearing Aids, \acrshort{HAs}) proposem una metodologia d'avaluació que, enlloc de senyals de test al laboratori, utilitza tècniques d'àudio espacial per generar escenes 3D realistes de parla amb soroll, les quals reproduïm amb un conjunt d'altaveus i re-enregistrem amb un cap de maniquí KU100, tant portant com sense portar audiòfons. Comparant els algorismes tradicionals dels audiòfons amb models prototip basats en DNN, mostrem que aquests últims generalitzen millor.

\item\textbf{\textit{(ii)}} Per a la millora de la parla monoaural per oients normals, examinem tres estratègies per augmentar la validesa ecològica de conjunts sintètics de respostes impulsionals de sala (Room Impulse Responses, \acrshort{RIR}s) —absorció multibanda, directivitat de la font i directivitat del receptor— i mostrem que emprar coeficients d'absorció dependents de la freqüència proporciona guanys mesurables quan s'avalua amb RIRs reals.

\item \textbf{\textit{(iii)}} Per a la separació de fonts musicals (Musical Source Separation, \acrshort{MSS}) d'enregistraments en directe des del públic, on els models entrenats amb gravacions d'estudi es degraden al ignorar l'acústica del recinte, la resposta del sistema de d'altaveus i el soroll del públic, introduïm CrowdioSet (pistes d'ambients reals combinats amb cors sintètics, generats a partir de conversions de veu cantada) i PaRIRset (respostes impulsionals estèreo captades en 40 recintes de concerts professionals reals), i mostrem que entrenar amb aquestes dades millora la separació en aquest cas d'ús.

\end{itemize}
En cadascuna de les tasques, el nostre enfoc es mostra efectiu i trasllada la complexitat de la inferència a la generació de les dades.